% Subsection: Comparing Closure Conditions  (notes stub; content authored later)
\subsection{Comparing Closure Conditions}

\begin{remark*}[Exposition]
Closure conditions are ordered by strength.  A stronger condition supplies a
larger range of allowed operations; a weaker condition supplies less internal
stability.  The purpose of comparing closure conditions is to know which
operations are legitimate inside a given set system.
\end{remark*}

\begin{remark*}[Comparison Table]
\begin{center}
\small
\renewcommand{\arraystretch}{1.2}
\begin{tabularx}{\linewidth}{>{\raggedright\arraybackslash}p{0.34\linewidth}
                            >{\raggedright\arraybackslash}X}
\toprule
\textbf{Closure condition} & \textbf{What it permits} \\
\midrule
Finite unions & Combine finitely many admitted alternatives. \\
Finite intersections & Combine finitely many admitted simultaneous conditions. \\
Complements & Replace an admitted condition by its ambient negation. \\
Countable unions & Combine a sequence of admitted alternatives. \\
Countable intersections & Combine a sequence of admitted simultaneous conditions. \\
Arbitrary unions & Combine any indexed collection of admitted alternatives. \\
\bottomrule
\end{tabularx}
\end{center}
\end{remark*}

\begin{remark*}[Implications]
Arbitrary union closure implies countable union closure, and countable union
closure implies finite union closure when the relevant finite family is viewed
as a countable family by repeating harmless terms or by using the finite case
directly.  The analogous implication holds from countable intersection closure
to finite intersection closure.
\end{remark*}

\begin{remark*}[Non-Implications]
Finite union closure alone does not imply complement closure.  Complement
closure alone does not imply union closure.  Countable union closure does not
imply arbitrary union closure.  These failures are why different theories
choose different closure packages.
\end{remark*}

\begin{remark*}[Forward Use]
The following sections package these closure conditions into named systems.
Semirings, rings, algebras of sets, sigma-algebras, and topological
structures differ mainly by which closure operations they require and which
generating procedures they allow.
\end{remark*}
